% P03.1 finite-flat local PL pullback proof, live integration version.

\begin{proof}[\Cref{lem:p031-characteristic-spectral-maximum}]
Fix $y\in U$, put $H:=\mathcal H(y)$, and base change
\eqref{eq:p031-characteristic-polynomial} to the finite $H$-algebra of the
analytic fiber.  After extension to an algebraic closure of $H$, its roots,
counted with multiplicity, are the residue-field images of $h$ at the points
of $\pi^{-1}(y)$.  Indeed, the fiber algebra is the product of its local
Artinian factors.  On each factor multiplication by $h$ is triangular after
scalar extension, with diagonal entries given by the residue-field
conjugates of $h$.  Nilpotents change multiplicities but introduce no new
eigenvalues.  Uniqueness of the extended absolute value on every finite
extension of the complete field $H$ shows that the largest root modulus is
\[
                       R_y:=\max_{v\in\pi^{-1}(y)}|h(v)|.
\]

Write the roots as $\alpha_1,\ldots,\alpha_d$.  Since $c_i(y)$ is, up to
sign, their $i$-th elementary symmetric function, the ultrametric inequality
gives
\[
                           |c_i(y)|\leq R_y^i.
\]
Conversely, choose a root $\alpha$ with $|\alpha|=R_y$.  If
$|c_i(y)|^{1/i}<R_y$ for every $i$, then in
\[
       \alpha^d+c_1(y)\alpha^{d-1}+\cdots+c_d(y)=0
\]
the first term would have strictly larger absolute value than every other
term, contradicting the ultrametric triangle law.  This proves
\eqref{eq:p031-characteristic-spectral-maximum}.

Multiplication by the unit $h$ is an automorphism of the finite locally free
module $\pi_*\mathcal O_V$, so $c_d$ is a unit.  Fix $y\in U$ and let
\[
                         I_y:=\{i:c_i(y)\ne0\}.
\]
After shrinking around $y$, each $c_i$ with $i\in I_y$ is a unit.  For
$i\notin I_y$, continuity and the nonvanishing of $c_d$ give
\[
                 |c_i(z)|^{1/i}<|c_d(z)|^{1/d}
\]
on a smaller neighborhood.  These indices do not contribute to the
spectral maximum there.  With $N:=\operatorname{lcm}(1,\ldots,d)$,
\begin{equation}
\label{eq:p031-local-characteristic-lfs-formula}
 \log R_z
  =\max_{i\in I_y}\frac1i\log|c_i(z)|
  =\frac1N\max_{i\in I_y}
       \{-\log|c_i(z)^{-N/i}|\}.
\end{equation}
This is a $\mathbf Q$-LFS presentation, proving the final assertion.
\end{proof}

\begin{proof}[\Cref{thm:p031-rational-lfs-fiber-maximum}]
Fix $x\in X$ and write $\pi^{-1}(x)=\{e_1,\ldots,e_q\}$.  Choose pairwise
disjoint ordinary open neighborhoods $V_k$ of the $e_k$ on which $u$ has
a presentation \eqref{eq:p031-rational-lfs-normal-form}.  Since a finite
morphism is proper,
\[
 U_0:=X\setminus\pi\bigl(E\setminus(V_1\cup\cdots\cup V_q)\bigr)
\]
is an open neighborhood of $x$, and $\pi^{-1}(U_0)$ is the disjoint union
of the $V_k\cap\pi^{-1}(U_0)$.

A finite flat analytic morphism is open by
\cite[proof of Proposition~5.11, p.~175]{GJR26}.  Shrink $U_0$ to the
intersection of the images of these finitely many branches and then along
their locally constant rank strata.  We obtain an open neighborhood $U$ of
$x$ and a decomposition
\begin{equation}
\label{eq:p031-local-finite-flat-branch-decomposition}
                 \pi^{-1}(U)=E_1\amalg\cdots\amalg E_q
\end{equation}
such that each $E_k\to U$ is finite locally free, surjective, and of a
fixed positive degree $d_k$.

On $E_k$, take a common denominator in the selected rational presentation
and distribute the finite maxima.  This gives
\begin{equation}
\label{eq:p031-rational-lfs-block-compression}
 u=\frac1{D_k}\max_{F\in\mathcal F_k}\{-\log|F|\},
 \qquad D_k\in\mathbf Z_{>0},
\end{equation}
where $\mathcal F_k$ is a finite family of analytic units on $E_k$.  For
the empty sum, take $\mathcal F_k=\{1\}$.  Applying
\Cref{lem:p031-characteristic-spectral-maximum} to $F^{-1}$ on every
branch gives
\begin{equation}
\label{eq:p031-branchwise-spectral-compression}
 \pi_{\max}u(z)
  =\max_{1\leq k\leq q}\max_{F\in\mathcal F_k}
      \frac1{D_k}
      \log\max_{e\in(E_k)_z}|F(e)^{-1}|.
\end{equation}
Each function inside the finite maxima is $\mathbf Q$-LFS.  A finite
maximum of rational presentations is again $\mathbf Q$-LFS: compress each
presentation as in \eqref{eq:p031-rational-lfs-block-compression}, take a
common denominator, and raise the finitely many analytic units to the
corresponding positive integral powers.  Hence
\eqref{eq:p031-branchwise-spectral-compression} is $\mathbf Q$-LFS on
$U$.  Since $x$ was arbitrary, the fiber-maximum assertion follows.

Pullback sends a base $\mathbf Q$-LFS minorant of $f$ to a cover minorant,
so $s_{\mathbf Q,X}\circ\pi\leq s_{\mathbf Q,E}$.  Conversely, if
$u\leq f\circ\pi$ is a cover competitor, then $\pi_{\max}u$ is a base
$\mathbf Q$-LFS competitor and
\[
 u(e)\leq\pi_{\max}u(\pi(e))
       \leq s_{\mathbf Q,X}(\pi(e)).
\]
Taking the supremum over $u$ proves the reverse inequality in
\eqref{eq:p031-rational-lfs-envelope-pullback}.  It also proves that the
two competitor families are empty simultaneously.
\end{proof}

\begin{proof}[\Cref{prop:p031-local-rational-lfs-cofinality}]
By \cite[\S7.9, p.~190]{GJR26}, after shrinking around $y$ one may write
\[
 u=\sum_{\nu=1}^{r}a_\nu\phi_\nu,
 \qquad
 \phi_\nu=\max_{1\leq j\leq N_\nu}\{-\log|f_{\nu j}|\},
 \qquad a_\nu\in\mathbf R_{\geq0},
\]
with analytic units $f_{\nu j}$.  Choose an arbitrarily small compact
strictly analytic neighborhood $U$ of $y$ inside this presentation domain.
The finitely many functions $\phi_\nu$ are bounded on $U$.

Put $\eta_n:=1/(4n)$ and choose $a_{\nu,n}\in\mathbf Q_{\geq0}$ so that
\[
 \left\|\sum_{\nu=1}^{r}(a_{\nu,n}-a_\nu)\phi_\nu
 \right\|_{U,\infty}<\eta_n.
\]
Set $w_n:=\sum_\nu a_{\nu,n}\phi_\nu$.  Nontriviality of the valuation
gives $\gamma\in K^\times$ with $\tau:=-\log|\gamma|>0$.  Choose
$q_n\in\mathbf Q_{>0}$ with
\[
                         \eta_n<q_n\tau<2\eta_n.
\]
The constant $-q_n\tau$ is $\mathbf Q$-LFS because
\[
              -q_n\tau=q_n\max\{-\log|\gamma^{-1}|\}.
\]
Thus $v_n:=w_n-q_n\tau$ is $\mathbf Q$-LFS and satisfies
$u-3\eta_n<v_n<u$ on $U$, which proves
\eqref{eq:p031-one-sided-local-rationalization}.

The rational local competitor family is contained in the LFS family.  For
the reverse inequality, apply the ordered approximation to an arbitrary
local LFS minorant of $g$.  Restrict $v_n$ to an ordinary open neighborhood
of $y$ contained in the interior of $U$.  It remains below that minorant and
hence below $g$, while $v_n(y)\to u(y)$.  Taking the supremum proves
\eqref{eq:p031-local-real-rational-lfs-cofinality}.  If the LFS family is
empty, so is its rational subfamily, proving the final clause.
\end{proof}

\begin{proof}[\Cref{thm:p031-finite-flat-local-envelope-pullback}]
Pullback of local LFS minorants gives
\begin{equation}
\label{eq:p031-finite-flat-easy-local-inequality}
 P_{\mathrm{LFS}}^{\mathrm{loc}}f(x)
   \leq P_{\mathrm{LFS}}^{\mathrm{loc}}(f\circ\pi)(e).
\end{equation}

For the reverse inequality, let $u$ be an LFS minorant of $f\circ\pi$ on
an ordinary open neighborhood of $e$.  Write
$\pi^{-1}(x)=\{e_1,\ldots,e_q\}$ with $e_1=e$, and let $b_0$ be the
assumed LFS minorant on a neighborhood of the complete fiber.  Choose
pairwise disjoint ordinary open neighborhoods $W_k$ of the $e_k$ such that
$u$ is defined on $W_1$ and $b_0$ is defined on every $W_k$ for $k>1$.
On their disjoint union, set
\[
 \widetilde u|_{W_1}:=u,
 \qquad
 \widetilde u|_{W_k}:=b_0\quad(k>1).
\]
This is an LFS minorant of $f\circ\pi$ near the complete fiber.

Fix $\varepsilon>0$.  Apply
\Cref{prop:p031-local-rational-lfs-cofinality} at each $e_k$, and restrict
the resulting compact neighborhoods to pairwise disjoint ordinary open
neighborhoods.  Their union is a neighborhood $W'$ of the complete fiber
carrying a $\mathbf Q$-LFS function $v$ such that
\begin{equation}
\label{eq:p031-complete-fiber-rational-minorant}
             v\leq\widetilde u\leq f\circ\pi,
             \qquad v(e)>u(e)-\varepsilon.
\end{equation}
Since $\pi$ is finite and hence proper, after shrinking to an ordinary open
neighborhood $U$ of $x$ one has $\pi^{-1}(U)\subset W'$.  By
\Cref{thm:p031-rational-lfs-fiber-maximum},
$m:=\pi_{\max}(v|_{\pi^{-1}(U)})$ is $\mathbf Q$-LFS on $U$, and
\[
                         m\leq f|_U,
             \qquad m(x)\geq v(e)>u(e)-\varepsilon.
\]
Therefore $u(e)-\varepsilon<P_{\mathrm{LFS}}^{\mathrm{loc}}f(x)$.
Letting $\varepsilon\downarrow0$ and taking the supremum over $u$ proves
the reverse of \eqref{eq:p031-finite-flat-easy-local-inequality}.

If both spaces are boundaryless, every point is interior.  Apply
\Cref{thm:p031-local-pl-lfs-cofinality} on $E$ and $X$ to the two sides of
\eqref{eq:p031-finite-flat-local-lfs-pullback}; this gives
\eqref{eq:p031-finite-flat-local-pl-pullback}.
\end{proof}

\begin{proof}[\Cref{thm:p031-finite-flat-pl-defect-balance}]
A global LFS minorant pulls back to a global cover LFS minorant, so
\Cref{thm:p031-finite-flat-local-envelope-pullback} applies at every cover
point.  It also makes all local and global envelope values finite.  Pullback
of global PL competitors gives
\[
                   s_{\mathrm{PL},X}\circ\pi
                      \leq s_{\mathrm{PL},E},
\]
so $\Gamma^{\mathrm{PL}}_{\pi,f}\geq0$.

Fix $e\in E$, put $x:=\pi(e)$, and denote the common value in
\eqref{eq:p031-finite-flat-local-pl-pullback} by $P_x$.  Then
\begin{align*}
 \Delta^{\mathrm{PL}}_{f,X}(x)
   -\Delta^{\mathrm{PL}}_{f\circ\pi,E}(e)
  &=(P_x-s_{\mathrm{PL},X}(x))
      -(P_x-s_{\mathrm{PL},E}(e))\\
  &=s_{\mathrm{PL},E}(e)-s_{\mathrm{PL},X}(x)\\
  &=\Gamma^{\mathrm{PL}}_{\pi,f}(e),
\end{align*}
which proves \eqref{eq:p031-finite-flat-pl-defect-balance}.  Since the
cover defect and the gain are nonnegative, the same identity gives
\eqref{eq:p031-finite-flat-pl-defect-contraction} and
\eqref{eq:p031-finite-flat-pl-gain-bound}.

Every LFS function is semipositive real PL, so the full P03 union envelope
is the PL envelope.  Its pullback defect is therefore the displayed gain,
which vanishes exactly when
\eqref{eq:p031-finite-flat-pl-defect-preservation} holds.  If the base
defect vanishes at $x$, the two nonnegative summands
\[
 \Delta^{\mathrm{PL}}_{f\circ\pi,E}(e)
       +\Gamma^{\mathrm{PL}}_{\pi,f}(e)
       =\Delta^{\mathrm{PL}}_{f,X}(x)
\]
vanish for every $e\in\pi^{-1}(x)$.
\end{proof}

\begin{proof}[\Cref{cor:p031-finite-orbit-four-defect-balance}]
For a finite analytic orbit presentation, the global PL pullback gain is
$\Sigma_{\pi,f}$ by \eqref{eq:p031-full-union-symmetry-defect}.  Substitute
the split identity \eqref{eq:p031-pl-lfs-defect-split} on $X$ and $E$ into
\eqref{eq:p031-finite-flat-pl-defect-balance}.  This gives
\eqref{eq:p031-finite-orbit-four-defect-balance} at every cover point.

If the base LFS globalization defect vanishes at $x$, then
\Cref{thm:p031-pl-lfs-defect-split} gives
\[
             \Delta^{\mathrm{PL}}_{f,X}(x)=B_{f,X}(x)=0.
\]
The finite-flat defect balance makes the cover PL defect and
$\Sigma_{\pi,f}$ vanish throughout the fiber.  The envelope chain
\eqref{eq:p031-lfs-pl-envelope-chain}, together with the zero base branch
gap and zero symmetry defect, then forces the cover branch gap to vanish.
The split identity makes the cover LFS defect vanish as well.  Exact
full-envelope pullback compatibility follows from
\eqref{eq:p031-full-union-symmetry-defect}.
\end{proof}
